\documentclass[10pt,a4paper]{amsart}
\usepackage[margin=19mm]{geometry}
\usepackage{amsmath,amssymb,mathtools}
\usepackage[scr=rsfs,cal=euler]{mathalfa}
\usepackage{xfrac}
\usepackage[unicode,hidelinks,bookmarks=false]{hyperref}
\newcommand{\ie}{i.e.\ }
\newcommand{\cf}{cf.\ }
\newcommand{\resp}{respectively\ }
\newcommand{\Subsection}{\subsection}
\providecommand{\nobreakdash}{\nobreak\hbox{-}}
\setlength{\emergencystretch}{2em}
\multlinegap0pt
\usepackage{amssymb}
\usepackage{mathtools}
\usepackage{xcolor}
\usepackage{enumerate}
\newtheorem{theo}{Theorem}
\newtheorem{lemm}{Lemma}
\newtheorem{prop}{Proposition}
\DeclareMathOperator{\Aut}{Aut}
\DeclareMathOperator{\IC}{IC}
\DeclareMathOperator{\IH}{IH}
\DeclareMathOperator{\MHM}{MHM}
\newcommand{\PP}{\mathbb{P}}
\newcommand{\QQ}{\mathbf{Q}}
\newcommand{\ZZ}{\mathbf{Z}}
\newcommand{\bft}{\mathbf{t}}
\newcommand{\cH}{\mathcal{H}}
\newcommand{\cK}{\mathcal{K}}
\newcommand{\cM}{\mathcal{M}}
\DeclareMathOperator{\gr}{gr}
\newcommand{\lround}[1]{\left\lfloor #1 \right\rfloor}
\newcommand{\lsta}{_{*}}
\newcommand{\norm}[1]{\|#1\|}
\DeclareMathOperator{\prim}{prim}
\newcommand{\uind}[1]{^{(#1)}}
\newcommand{\uround}[1]{\left\lceil #1 \right\rceil }
\title{Hodge theory and $K$-stability of some very symmetric hypersurfaces}
\author{Hyunsuk Kim}
\date{Source: arXiv:2604.27229v1 (CC BY 4.0)}
\begin{document}
\maketitle

\noindent\emph{Source and licence.} Hodge theory and $K$-stability of some very symmetric hypersurfaces. Version arXiv:2604.27229v1 (CC BY 4.0), distributed by arXiv under the Creative Commons Attribution 4.0 International licence, http://creativecommons.org/licenses/by/4.0/, as recorded in the arXiv licence metadata for this exact version and shown on the abstract page https://arxiv.org/abs/2604.27229v1. Full licence notice: \texttt{licenses/arxiv-2604.27229-CC-BY-4.0.txt}.\par\smallskip

\noindent\textit{Editorial scope.} Counted unit: the article's prop prop:intersection-cohomology (source0.tex lines 884--892). The article's complete original proof of it is delivered with it (source0.tex lines 914--926, one \texttt{proof} environment of 13 lines, which the article places away from the statement). No mathematics is written, repaired or re-derived: the statement and the proof are the article's own lines, in the article's own notation, and no retained line is substituted or re-flowed.  Retained uncounted as the source-local context the proof needs: lines 411--421; lines 440--450; lines 826--829; lines 833--842; lines 870--874; lines 883--910; lines 937--939.  Nothing was omitted from any retained span, so the package carries no omission record.  No retained line contains a citation, so the wrapper carries no bibliography and no bibliography entry.  No retained line contains a figure environment, an image-inclusion command or drawing code, so no image is delivered and the package declares no asset dependency.  Editorial formatting only, in addition: an AMS \texttt{amsart} preamble that restates the article's own declarations and macros used by the retained lines, namely the environments \texttt{theo}, \texttt{lemm}, \texttt{prop} on separate counters, together with the standard packages those lines need.  The retained environments print in this document as Theo 1, Lemm 1, Proposition 1 in order of appearance, whereas the article prints the same items with the numbering of its own full document.  The complete native source of the article as distributed in the arXiv e-print for this exact version is bundled unchanged as \texttt{sources/199457/source0.tex}.
    \begin{theo} \label{theo:appending-snc}
        Let $X \subset \PP^{m+1}$ be a hypersurface given by the homogeneous polynomial $f(x)$ of degree $d$. Consider the hypersurface
        $$ Y = \{ f(x) + y_{11} \cdots y_{1d} + \ldots + y_{l1}\cdots y_{ld} = 0\} \subset \PP^{m+ 1+ld}.$$
        The singular cohomology of $X$ and $Y$ is related in the following manner:
        $$ \overline{H}^{m+ld+ i}(Y) \simeq \bigoplus_{0 \leq i_{0}} \overline{H}^{m + i_{0}}(X) (-l-i+i_{0})^{\oplus a_{i-i_{0}}}.$$
        In particular, if $X$ is smooth, we have
        $$ \overline{H}^{m+ld + i}(Y) \simeq H^{m}_{\prim}(X)(-l-i)^{\oplus a_{i}}.$$
        Here, the multiplicities $a_{i}$ are defined as
        $$ a_{i} = \sum_{\substack{i_{1}+ \ldots + i_{l} = i \\ 0 \leq i_{j} \leq d-2}} {d - 1 \choose i_{1}}\cdots {d - 1 \choose i_{l}}.$$
        These numbers are non-zero for $0 \leq i \leq l(d-2)$.
    \end{theo}

    \begin{theo}\label{theo:main-theorem}
        Let $X_{l,d}$ be the hypersurface in $\PP^{ld-1}$ defined by the equation
        $$ x_{11} \cdots x_{1d} + \ldots + x_{l1}\cdots x_{ld} = 0.$$
        \begin{enumerate}
            \item The automorphism group $\Aut(X_{l,d})$ is a semi-direct product of a finite group and a torus. For the exact description of the group, see Proposition \ref{prop:automorphism-of-X}.
            \item The $\QQ$-Hodge structure of the singular cohomology of $X_{l,d}$ is given by
            $$ \overline{H}^{ld-2+i}(X_{l,d}, \QQ) \simeq \QQ(-(l-1+i))^{\oplus a_{i}}$$
            where $a_{i}$ is defined as in Theorem \ref{theo:appending-snc}.
            \item The intersection cohomology of $X_{l,d}$ is pure of Hodge--Tate type, and we have an explicit formula for the intersection Betti numbers.  We refer to Proposition \ref{prop:intersection-cohomology} for the formula.
        \end{enumerate}
    \end{theo}

     We give a small remark which follows from standard $\mathfrak{sl}_{2}$-representation theory.
     \begin{lemm} \label{lemm:inj-surj-of-V-uind(m)}
         $V_{\bft}\uind{m} \to V_{\bft}\uind{m-1}$ is surjective if and only if $m \leq \uround{\frac{|\bft|}{2}} - l$ and is injective if and only if $m \geq \lround{\frac{|\bft|}{2}} - l + 1$. In particular, $V_{\bft,\prim}\uind{m}$ is non-zero for $0 \leq m \leq \lround{\frac{|\bft|}{2}}-l$ where $|\bft|= t_{1}+ \ldots + t_{l}$.
     \end{lemm}

     \begin{lemm} \label{lemm:dimension-of-Vt-prim}
         $$ \sum_{m} \dim V_{\bft, \prim}\uind{m} q^m = \tau_{<\frac{|\bft|}{2}-l+1} \left( (1-q) \prod_{i=1}^{l} \frac{q^{t_{i}-1}-1}{q-1} \right) \in \ZZ[q].$$
         The truncation $\tau_{<\frac{|\bft|}{2}-l+1}$ of a polynomial is defined as
         $$ \tau_{<\frac{|\bft|}{2}-l+1} \left( \sum_{m \geq 0} \alpha_{m} q^{m} \right) = \sum_{m < \frac{|\bft|}{2}-l+1} \alpha_{m}q^{m}.$$
     \end{lemm}
     \begin{proof}
         It is easy to see that
         $$ \sum_{m} \dim V_{\bft}\uind{m} q^{m} = \prod_{i=1}^{l} \frac{q^{t_{i}-1}-1}{q-1} \in \ZZ[q].$$
         The assertion then follows from Lemma \ref{lemm:inj-surj-of-V-uind(m)}.
     \end{proof}

    \begin{prop} \label{prop:RHM-defect-object-X-graded}
        We have the following description of the RHM defect object $\cK_{X}$ of $X$:
        $$ \bigoplus_{w} \gr_{w}^{W} \cK_{X} \simeq \bigoplus_{\bft \neq (d, \ldots, d)} \bigoplus_{|J_{i}| = t_{i}} \bigoplus_{m \geq 0} \IC_{\PP(L_{J})}^{H} (-m-l+1) \otimes_{\QQ} V_{\bft, \prim}\uind{m}.$$
        Here, $\PP(L_{J})$ is the projectivization of $L_{J}$.
    \end{prop}

    \subsection{Intersection cohomology}
    \begin{prop} \label{prop:intersection-cohomology}
    The intersection cohomology $\IH^{\bullet}(X)$ is of Hodge--Tate type. Furthermore, we have the following formula for the generating series of the intersection Betti numbers.
    $$
        \sum \dim \IH^{2i}(X) q^{i} =\frac{q^{ld-1}-1}{q - 1}+ (q^{d-1} - (q-1)^{d-1})^{l} q^{l-1} + \sum_{\nu} {l \choose \nu_2, \nu_3, \ldots, \nu_d } {d \choose 2}^{\nu_{2}} \cdots {d \choose d}^{\nu_{d}} g_{\nu}(q) \in \ZZ[q],
    $$
    where
    $$ g_{\nu}(q) = (-1)^{\norm{\nu}} \frac{q^{ld -\norm{\nu}}-1}{q-1} \cdot \tau_{< \norm{\nu}/2} \left( (1-q)\cdot  q^{l-1} \prod_{k=2}^{d} \left(\frac{q^{k-1}-1}{q-1}\right)^{\nu_{k}}  \right).$$
    and the sum ranges over $\nu = (\nu_{2},\ldots, \nu_{d}) \in \ZZ_{\geq 0}^{d-1}$ such that $\nu_{2} + \ldots + \nu_{d} = l$ and $\nu_{d} \neq l$. Here, $\norm{\nu}  = \sum_{k=2}^{d} k \nu_{k}$.
    \end{prop}
    We first compute the contribution of the defect object in the Grothendieck group of mixed Hodge structures $K_{0}(\mathrm{MHS})$.
    \begin{lemm} \label{lemm:defect-object-Grothendieck-group}
        We have
        $$ (-1)^{ld-1} \sum_{i} (-1)^{i} [\cH^{i} a\lsta \cK_{X}] = \sum_{\nu} {l \choose \nu_{2},\ldots, \nu_{d}} {d \choose 2}^{\nu_{2}} \cdots {d \choose d}^{\nu_{d}} g_{\nu}(u) \in K_{0}(\mathrm{MHS}),$$
        where $g_{\nu}(u)$ is defined as in Proposition \ref{prop:intersection-cohomology}. Here, $u = \QQ(-1) \in K_{0}(\mathrm{MHS})$. The sum ranges over $\nu = (\nu_{2},\ldots, \nu_{d}) \in \ZZ_{\geq 0}^{d-1}$ such that $\nu_{2} + \ldots + \nu_{d} = l$ and $\nu_{d} \neq l$.
    \end{lemm}
    \begin{proof}
        The alternating sum of cohomology gives a homomorphism 
        $$\chi \colon K_{0}(\MHM(X)) \to K_{0}(\mathrm{MHS}), \qquad [\cM] \mapsto \sum (-1)^{i} [\cH^{i} a\lsta \cM ]$$
        from the Grothendieck group of mixed Hodge modules on $X$ to the Grothendieck group of mixed Hodge structures. Therefore, it is enough to calculate the image of $\bigoplus_{w} \gr_{w} \cK_{X}$ under this homomorphism. From Proposition \ref{prop:RHM-defect-object-X-graded}, we have
        $$ \chi (\cK_{X}) = \sum_{\bft \neq (d,\ldots, d)} \bigoplus_{|J_{i}| = t_{i}} \bigoplus_{m\geq 0} \chi\left(\IC_{\PP(L_{J})}^{H}(-m-l+1)\right)\cdot \dim V_{\bft, \prim}\uind{m}.$$
        Note that $\dim L_{J} = ld - |J|$, therefore $$\chi\left(\IC_{\PP(L_{J})}^{H}(-m-l+1)\right) = (-1)^{ld- |J|-1}u^{m+l-1}\frac{u^{ld -|J|}-1}{u-1}.$$
        Let $\nu_{k} = \#\{ i : t_{i} = k\}$. For $\nu = (\nu_{2},\ldots, \nu_{d})$ such that $\nu_{2} + \ldots + \nu_{d} = l$, there are
        $${l \choose \nu_{2},\ldots, \nu_{d}} {d \choose 2}^{\nu_{2}} \cdots {d \choose d}^{\nu_{d}}$$
        many possibilities for $J = (J_{1},\ldots, J_{l})$ having such value of $\nu$. Then the assertion follows by applying Lemma \ref{lemm:dimension-of-Vt-prim} and noting that
        $$ \sum_{m\geq 0} u^{m+l-1} \dim V_{\bft, \prim}\uind{m} = \tau_{<\frac{|\bft|}{2}} \left( u^{l-1} (1-u) \prod_{i=1}^{l} \frac{u^{t_{i}-1}-1}{u-1} \right) = \tau_{<\frac{\norm{\nu}}{2}} \left( u^{l-1} (1-u) \prod_{k=2}^{d} \left( \frac{u^{k-1}-1}{u-1} \right)^{\nu_{k}} \right), $$
        and $|J| = |\bft| = \norm{\nu}$.
    \end{proof}

    \begin{prop} \label{prop:automorphism-of-X}
        Let $G$ and $X$ as above. Then $\Aut(X) = G$.
    \end{prop}

    \begin{prop} 
    The intersection cohomology $\IH^{\bullet}(X)$ is of Hodge--Tate type. Furthermore, we have the following formula for the generating series of the intersection Betti numbers.
    $$
        \sum \dim \IH^{2i}(X) q^{i} =\frac{q^{ld-1}-1}{q - 1}+ (q^{d-1} - (q-1)^{d-1})^{l} q^{l-1} + \sum_{\nu} {l \choose \nu_2, \nu_3, \ldots, \nu_d } {d \choose 2}^{\nu_{2}} \cdots {d \choose d}^{\nu_{d}} g_{\nu}(q) \in \ZZ[q],
    $$
    where
    $$ g_{\nu}(q) = (-1)^{\norm{\nu}} \frac{q^{ld -\norm{\nu}}-1}{q-1} \cdot \tau_{< \norm{\nu}/2} \left( (1-q)\cdot  q^{l-1} \prod_{k=2}^{d} \left(\frac{q^{k-1}-1}{q-1}\right)^{\nu_{k}}  \right).$$
    and the sum ranges over $\nu = (\nu_{2},\ldots, \nu_{d}) \in \ZZ_{\geq 0}^{d-1}$ such that $\nu_{2} + \ldots + \nu_{d} = l$ and $\nu_{d} \neq l$. Here, $\norm{\nu}  = \sum_{k=2}^{d} k \nu_{k}$.
    \end{prop}

    \begin{proof}[Proof of Proposition \ref{prop:intersection-cohomology}]
        From the short exact sequence
        $$ 0 \to \cK_{X} \to \QQ_{X}^{H}[ld-2] \to \IC_{X}^{H} \to 0,$$
        we have
        $$ \sum_{i} (-1)^{i} [\IH^{i}(X)] = (-1)^{ld- 1} \sum_{i} (-1)^{i} [\cH^{i} a\lsta \cK_{X}] + \sum_{i} (-1)^{i} [H^{i}(X)] \in K_{0}(\mathrm{MHS}).$$
        The first object is calculated in Lemma \ref{lemm:defect-object-Grothendieck-group}. We compute the second term using Theorem \ref{theo:main-theorem}(2). We have
        \begin{align*}
            &\sum_{i} (-1)^{i} [H^{i}(X)] \\
            & = \frac{u^{ld-1}-1}{u-1} + (-1)^{ld-2} u^{l-1} (a_{0} - a_{1}u + \ldots + (-1)^{l(d-2)} a_{l(d-2)}u^{l(d-2)}) \\
            & = \frac{u^{ld-1}-1}{u-1} + (-1)^{ld-2} u^{l-1} \left( {d-1 \choose 0} - {d-1 \choose 1} u + \ldots + (-1)^{d-2} {d-1 \choose d-2} u^{d-2} \right)^{l},
        \end{align*}
        where the second equality follows from the definition of $a_{i}$ in Theorem \ref{theo:appending-snc}. The rest of the proof is a straightforward computation.
    \end{proof}

\end{document}
